\documentclass[11pt]{article}
\usepackage[a4paper,margin=1in]{geometry}
\usepackage[T1]{fontenc}
\usepackage{lmodern}
\usepackage{amsmath,amssymb,amsthm,mathtools}
\usepackage{microtype}
\usepackage{enumitem}
\usepackage{xcolor}
\usepackage[numbers,sort&compress]{natbib}
\usepackage[colorlinks=true,linkcolor=blue!55!black,citecolor=blue!55!black,urlcolor=blue!55!black]{hyperref}
\usepackage[nameinlink,noabbrev]{cleveref}

\newtheorem{theorem}{Theorem}[section]
\newtheorem{proposition}[theorem]{Proposition}
\newtheorem{corollary}[theorem]{Corollary}
\theoremstyle{definition}
\newtheorem{definition}[theorem]{Definition}
\newtheorem{remark}[theorem]{Remark}
\newcommand{\e}{\mathrm e}
\newcommand{\T}{\mathbb T}
\newcommand{\norm}[1]{\lVert#1\rVert}

\title{\textbf{A Parseval Corridor for Direct Additive Twists\\
on the Determinant-Two M\"obius Core}}
\author{Liang Wang\\
\small School of Artificial Intelligence and Automation,\\[-2pt]
\small Huazhong University of Science and Technology, Wuhan 430074, P.R. China\\
\small \texttt{wangliang.f@gmail.com}}
\date{July 2026}

\begin{document}
\maketitle

\begin{abstract}
The periodic-approximation route to an additive phase has a sharp
minor-arc obstruction.  We open a different, direct route by treating
the phase as an \(L^2\) variable.  On the literal determinant-two
two-M\"obius core, Parseval gives an exact mean-square identity for
the normalized additive twist, at every scale and without a
scale-exceptional set.  The same identity holds on every complete
Fourier grid at least as long as the fiber interval.  In the
small-polylogarithmic envelope this is a positive fixed-\(X\) power
in phase \(L^2\).  It is not a pointwise theorem for a specified
production phase, and therefore does not close the original
pointwise additive-twist frontier.
\end{abstract}

\noindent\textbf{Keywords.}
M\"obius correlation; additive twist; Parseval identity; phase
average; determinant two.

\medskip
\noindent\fcolorbox{red!70!black}{red!3}{\parbox{0.94\linewidth}{
\textbf{Claim boundary.}
The power saving below is in Lebesgue \(L^2\) over the additive phase.
It is not a bound at a distinguished phase and is not a production
phase-cell registry.  Parseval uses orthogonality rather than a new
pointwise M\"obius estimate.  No physical H3 return, deterministic
all-prefix theorem, \(1/400\), prime-pair lower bound, or twin-prime
theorem is proved.}}

\section{The literal core and its phase transform}

Put \(\e(x)=\exp(2\pi i x)\).  Retain the determinant-two data
\[
 a,s\ge1,\quad d,u\in\mathbb Z,\quad
 (a,s)=1,\quad as\ {\rm odd},\quad su-ad=2,
\]
\[
 q=as,\qquad t(z)=ad+qz,\qquad
 c_z=\mu(d+sz)\mu(u+az).
\]
For \(N>0\), let
\[
 I_N=\{z\in\mathbb Z:N<t(z)\le2N\},\quad
 L_N=|I_N|,\quad E_N=\sum_{z\in I_N}|c_z|^2.
\]
The interval \(I_N\) consists of consecutive integers, and
\begin{equation}
 E_N\le L_N\le\frac Nq+1.
\label{eq:count}
\end{equation}
Define the normalized direct additive twist
\begin{equation}
 F_N(\alpha)=\frac qN\sum_{z\in I_N}c_z\e(-\alpha z),
 \qquad \alpha\in\T=\mathbb R/\mathbb Z.
\label{eq:transform}
\end{equation}

TPC-158 controls explicit major arcs through periodic approximation
and proves that this approximation route is expensive on a uniform
minor-arc family \citep{WangTPC158}.  The definition
\eqref{eq:transform} does not approximate the phase.

\section{Exact phase mean square}

\begin{theorem}[Direct additive-twist Parseval identity]\label{thm:parseval}
For every \(N>0\),
\begin{equation}
 \boxed{\quad
 \int_{\T}|F_N(\alpha)|^2\,d\alpha
 =\frac{q^2}{N^2}E_N.
 \quad}
\label{eq:parseval}
\end{equation}
In particular,
\begin{equation}
 \norm{F_N}_{L^2(\T)}^2
 \le\frac qN+\frac{q^2}{N^2}.
\label{eq:l2bound}
\end{equation}
No exceptional set of scales is required.
\end{theorem}

\begin{proof}
Expanding the square and using character orthogonality gives
\[
 \int_{\T}\e\bigl(-\alpha(z-z')\bigr)\,d\alpha
 =\mathbf1_{\{z=z'\}}.
\]
Thus every off-diagonal term vanishes and the diagonal equals
\(q^2E_N/N^2\).  Apply \eqref{eq:count}.
\end{proof}

\begin{corollary}[Exceptional phases]\label{cor:measure}
For every \(\lambda>0\),
\begin{equation}
 \operatorname{meas}\{\alpha\in\T:|F_N(\alpha)|>\lambda\}
 \le
 \frac{q^2E_N}{N^2\lambda^2}
 \le
 \frac{q/N+q^2/N^2}{\lambda^2}.
\label{eq:chebyshev}
\end{equation}
\end{corollary}

\begin{proof}
This is Chebyshev's inequality applied to \eqref{eq:parseval}.
\end{proof}

The identity is stronger than a logarithmic estimate in its own
declared norm, but weaker than a bound at any named \(\alpha\).
It holds equally for arbitrary bounded coefficients; the
M\"obius form identifies the actual core rather than supplying the
orthogonality.

\section{A finite exact Fourier registry}

Write \(I_N=\{z_0,\ldots,z_0+L_N-1\}\).

\begin{proposition}[Complete-grid identity]\label{prop:grid}
If \(M\ge L_N\), then
\begin{equation}
 \frac1M\sum_{r=0}^{M-1}
 \left|F_N\left(\frac rM\right)\right|^2
 =\frac{q^2}{N^2}E_N.
\label{eq:grid}
\end{equation}
Consequently the number of grid phases for which
\(|F_N(r/M)|>\lambda\) is at most
\[
 \frac{Mq^2E_N}{N^2\lambda^2}.
\]
\end{proposition}

\begin{proof}
After expansion, the finite character average is one precisely when
\(M\mid z-z'\).  Since \(|z-z'|\le L_N-1<M\), this occurs precisely
on the diagonal.  The counting assertion follows by summing
\(|F_N|^2>\lambda^2\) over bad grid points.
\end{proof}

The condition \(M\ge L_N\) is structural.  For a shorter grid,
distinct exponents may collide modulo \(M\), so \eqref{eq:grid} need
not hold.

\section{Interpretation of the power envelope}

Suppose the small-polylogarithmic envelope
\[
 q\le(\log X)^{\eta_0},\qquad N\ge\sqrt X
\]
holds.  For sufficiently large \(X\), \(q\le N\), and
\eqref{eq:l2bound} gives
\begin{equation}
 \boxed{\quad
 \norm{F_N}_{L^2(\T)}
 \le\sqrt2\,X^{-1/4}(\log X)^{\eta_0/2}.
 \quad}
\label{eq:power}
\end{equation}
Thus TPC-167 has the precise program status
\[
 \mathsf{PROVED\_L1\_ACTUAL\_CORE\_
 PHASE\_METRIC\_SINGLE\_CELL}.
\]
Its analytic norm is \(\mathsf{L2\_PHASE}\), while
\(\mathsf{program\_positive\_L2=false}\) and
\(\mathsf{fixed\_atom=false}\).  The positive exponent in
\eqref{eq:power} belongs to a phase-averaged norm.  It cannot be
entered as a production fixed-phase or endpoint-V3 payment.

\section{Why pointwise selection remains open}

The implication
\[
 \norm{F_N}_{L^2(\T)}\ll N^{-1/2}q^{1/2}
 \quad\Longrightarrow\quad
 |F_N(\alpha_\star)|=o(1)
\]
is false without information about the selected phase
\(\alpha_\star\).  At the coefficient level, take \(c_z=1\) on an
interval and \(\alpha_\star=0\).  Then
\[
 |F_N(0)|=\frac{qL_N}{N},
\]
which is of constant order when \(L_N\sim N/q\), while Parseval still
has mean square of order \(q/N\).  This example is a logical
nonimplication for coefficient-blind phase averaging; it is not a
claim that the literal M\"obius coefficients realize the example.

Accordingly, the parent-ready direct-twist node is advanced only to
a typed child:
\[
 \mathsf{direct\ additive\ twist}
 \longrightarrow
 \mathsf{phase\mbox{-}L^2\ direct\ twist}.
\]
The original pointwise node still needs a source-locked phase
registry plus an avoidance or selection theorem, or a genuinely
pointwise arithmetic estimate.

\section{Reproducible certificate}

The standard-library audit source-locks TPC-158, checks Parseval on a
shifted signed/zero fixture using two complete Fourier grids of
different sizes, and verifies the normalization bound.  These are
finite implementation checks; the proof above establishes the
continuous identity.  Its executable mutation contract rejects
promotion from phase \(L^2\) to a
fixed phase, from a phase power to endpoint V3, or from
orthogonality to M\"obius-specific cancellation.  Hashes certify
artifact integrity only.

\section{Conclusion}

The minor-arc obstruction to small-period approximation does not stop
every direct route.  Parseval gives an exact, exceptional-scale-free
direct additive-twist theorem on the actual core, with a positive
fixed-\(X\) power in phase \(L^2\).  The cost of this advance is
equally exact: a distinguished physical phase remains outside the
theorem.  TPC-168 asks how much of this phase-average information
survives on a finite source-locked phase registry.

\bibliographystyle{plainnat}
\bibliography{references}
\end{document}
